\documentclass[11pt,a4paper]{article}
\usepackage{gtu-el-yazisi}
\begin{document}
\notbilgi{Sayısal Haberleşme}{Düz tepeli örnekleme, PAM/PWM/PPM, TDM ve PCM}{Ders-5.pdf}
\textit{Not: Orijinal dosyadaki basılı sayfalar (s. 9) dijitalleştirilmedi; yalnızca el yazısı sayfalar aktarıldı.}

\sayfa{1}
\section{Sample and Hold operation (örnekle ve tut işlemi)}

\subsection{Düz Tepeli örnekleme (Flat top sampling) (FTS)}

\sekilnotu[1]{Blok diyagram: soldan giriş çizgisi, ortada ``FTS ($\tau$)'' kutusu, sağa çıkış çizgisi; kırmızı not ``Flattop sampling'' kutuya işaret ediyor. Kutunun üstünde yüksekliği 1, genişliği $\tau$ olan dikdörtgen darbe çizilmiş. Kutunun altında ``Clock \quad $f_s \geq 2 f_m$'' kutusu, yukarı okla FTS kutusuna bağlı.}

\sekilnotu[1]{$x(t)$ eğrisi ve üzerinde düz tepeli (dikdörtgen) darbeler, kırmızı; genlikleri $x(0)$, $x(T_s)$, $x(2T_s)$, $x(3T_s)$. Zaman ekseninde $0$, $\tau$, $T_s$, $T_s+\tau$, $2T_s$, $2T_s+\tau$, $3T_s$, $3T_s+\tau$ işaretli; $\tau$ aralığı çift yönlü okla gösterilmiş, altına ``$\tau$ genişlikli darbe'' yazılmış. Kırmızı not: $\ast$ tepeyi takip etmiyor.}

\sekilnotu[1]{Başlık: Sample and hold circuit (örnekle ve tut devresi). $x(t)$ eğrisi ve altında düz tepeli darbeler; eksende $0$, $\tau$, $T_s$, $2T_s$, $\dots$, $kT_s$ işaretli.}

$x_s(t)$: Darbe Genlik mod (DGM)\\
\hspace*{2em}Pulse amplitude mod

Burda darbeler düz tepeli

bu discrete time analog modulation \quad zaman ayrık ama genlik değil

Discrete time analog modulation

\subsection{Sample/ Hold (Dump) circuit}

\sekilnotu[1]{Kesikli çizgili kutu (FTS): solda $x(t)$ girişi, $t=nT_s$ anında kapanan anahtar; ardından toprağa bağlı kondansatör ve onunla paralel ikinci anahtar ($t=nT_s+\tau$ anında kapanıyor); çıkış $x_s(t)$. Notlar: ``$nT_s$ anında hemen kapanıp açılıyor'' (birinci anahtar), ``$\tau$ saniye sonra bu anahtar kapanıyor ve gerilim sıfır oluyor'' (ikinci anahtar).}

\sekilnotu[1]{Blok diyagram: $x(t)$ girişi çarpıcıya ($\otimes$) girer, çarpıcıya alttan $x_\delta(t)=\sum_n \delta(t-nT_s)$ uygulanır; çarpıcı çıkışı $g(t)$; ardından içinde yüksekliği 1, genişliği $\tau$ olan $q(t)$ darbesi bulunan blok; çıkış $x_s(t)$. Bloğun altında $q(t) \donusum Q(f)$ yazıyor.}

\begin{align*}
g(t) &= x(t)\, x_\delta(t) = \sum_n x(nT_s)\, \underline{\delta(t-nT_s)}\\
x_s(t) &= q(t) \circledast g(t) = \sum_n x(nT_s)\, \underline{q(t-nT_s)}
\end{align*}
\begin{align*}
x_s(t) &= q(t) \circledast \bigl[x(t)\, x_\delta(t)\bigr] = q(t) \circledast \Bigl[\sum_n x(nT_s)\,\delta(t-nT_s)\Bigr]\\
&= \sum_n x(nT_s)\, q(t-nT_s)
\end{align*}

\sayfa{2}
\begin{align*}
X_s(f) &= Q(f)\,\bigl[X(f) \circledast X_\delta(f)\bigr] = Q(f)\,\Bigl[X(f) \circledast f_s \sum_n \delta(f-nf_s)\Bigr]\\
X_s(f) &= Q(f)\,\Bigl[f_s \sum_n X(f-nf_s)\Bigr]
\end{align*}
\[ \hookrightarrow\ \tau\,\operatorname{sinc}(f\tau)\, e^{-j2\pi f \tau/2} \]
\[ X_s(f) = \underbrace{f_s\, Q(f)\, X(f)}_{!!} + Q(f)\cdot f_s \sum_{n\neq 0} X(f-nf_s) \]

\sekilnotu[2]{Frekans ekseninde dikdörtgen: $-f_m$ ile $f_m$ arasında yükseklik $T_s$; yanında ``alçak geçiren'' yazıyor.}

\vurgu[turuncu]{Equalizing LPF (Dengeleyici AGF)}
\[ x_s(t) \longrightarrow \underbrace{\boxed{H_{eq}(f)=\frac{1}{Q(f)}} \longrightarrow \boxed{\begin{matrix}\text{Ideal LPF}\\ \text{Gain}=T_s\\ \text{BW}=f_m\end{matrix}}}_{} \longrightarrow x(t) \]

\[ H_{eq}(f) = \begin{cases} \dfrac{T_s}{Q(f)} & |f|\le f_m\\[1ex] 0 & \text{else} \end{cases} \]

\sekilnotu[2]{$|Q(f)| = \tau\,|\operatorname{sinc}(f\tau)|$ grafiği: sinc genliği, ana lob tepesi $\tau$; sıfır geçişleri $-1/\tau$ ve $1/\tau$; eksende $-1/\tau$, $-f_m$, $f_m$, $1/\tau$ işaretli. $-f_m$ ve $f_m$ noktalarında kırmızı kesikli dikey çizgiler, ana lobun $|f|\le f_m$ aralığındaki tepe kısmı kırmızı kesikli çizilmiş ve kırmızı okla sağdaki $|H_{eq}(f)|$ grafiğine bağlanmış.}

\[ H_{eq}(f) = \begin{cases} \dfrac{T_s}{Q(f)} & |f|\le f_m\\[1ex] 0 & \text{else} \end{cases} \]

\sekilnotu[2]{$|H_{eq}(f)|$ grafiği: $-f_m$ ile $f_m$ arasında, ortası çukur dikdörtgen benzeri yanıt (merkezde minimum değer $T_s/\tau$, kenarlara doğru yükseliyor), dışarısı sıfır; yatay eksen $f$, işaretler $-f_m$ ve $f_m$.}

\[ H_{eq}(f) = \frac{T_s}{Q(f)} = \frac{T_s}{\tau\,\operatorname{sinc}(f\tau)}\,\underline{e^{\,j2\pi f \tau/2}}\,,\quad |f|\le f_m \]
\hspace*{3em}$\hookrightarrow$ advance in time $\rightarrow$ $h_{eq}(t)$ is not causal

For this reason, $H_{eq}'(f) = H_{eq}(f)\, e^{-j2\pi f T_x}$ is chosen\\
$\downarrow$
\[ h_{eq}'(t) = h_{eq}(t-T_x) \qquad T_x = \frac{\tau}{2}+T_0 \]
\[ x_s(t) \longrightarrow \boxed{H_{eq}'(f)} \longrightarrow x(t-T_0) \]

\sayfa{3}
\subsection{Reminder}
\sekilnotu[3]{Üçgen darbe: tepe yüksekliği $A$ ($t=0$'da), taban $-T$ ile $T$ arasında; yatay eksen $t$, işaretler $-T$, $0$, $T$.}
\[ A\,\Lambda\!\left(\frac{t}{T}\right) \overset{\mathcal{F}}{\longleftrightarrow} A\,T\,\operatorname{sinc}^2(fT) \]

\begin{ornek}[(2004 Midterm Ex.Q.)]
A PAM signal is produced as shown in the figure below. a) If the bandwidth of $x(t)$ is 1\,kHz, what is the maximum value of $T_s$ to obtain $x(t)$ without distortion at the receiver? b) For $T_s=T_s(\max)$, sketch PAM demodulator. Find the transfer function of demodulator.

\sekilnotu[3]{Blok diyagram: $x(t)$ girişi çarpıcıya ($\otimes$) girer, çarpıcıya alttan $\sum_n \delta(t-nT_s)$ uygulanır; çarpıcı çıkışı $q(t)$ bloğuna gider, çıkış $x_{PAM}(t)$. Sağda $q(t)$ grafiği: $t=0$'dan $0{,}2T_s$'e uzanan üçgen darbe, tepe noktası $t=0{,}1T_s$ anında ve yüksekliği $0{,}1T_s$; eksen işaretleri $0$, $0{,}1T_s$, $0{,}2T_s$.}
\end{ornek}

\begin{cozum}
\sekilnotu[3]{Aynı blok diyagram yeniden çizilmiş: $x(t) \to \otimes \to q(t) \to x_{PAM}(t)$, çarpıcıya alttan $\sum_n \delta(t-nT_s)$.}

a) $\mathrm{BW}\{x(t)\} = 1\ \mathrm{kHz} = f_m$\\
\hspace*{1.2em}$f_s \ge 2f_m = 2\ \mathrm{kHz}$\\
\hspace*{1.2em}$T_s \le \dfrac{1}{2f_m} = 0{,}5\ \mathrm{msec}$\\
\hspace*{1.2em}$T_s(\max) = 0{,}5\ \mathrm{msec}$

\sekilnotu[3]{$q(t)$ grafiği (sağda): $t=0$'dan $0{,}2T_s$'e kadar üçgen darbe, tepe yüksekliği $0{,}1T_s$ ($t=0{,}1T_s$ anında); eksen işaretleri $0$, $0{,}1T_s$, $0{,}2T_s$.}

b)
\[ x_{pam}(t) \longrightarrow \boxed{H_{eq}(f)=\frac{1}{Q(f)}} \longrightarrow \boxed{\begin{matrix}\text{Ideal}\\ \text{LPF}\\ \mathrm{BW}=f_m\end{matrix}} \longrightarrow x(t) \]
\[ H_{eq}(f) = \begin{cases} \dfrac{T_s}{Q(f)}, & |f| \le f_m\\[1ex] 0, & |f| > f_m \end{cases} \]
\begin{align*}
Q(f) &= 0{,}1T_s\,\Bigl(0{,}1T_s\operatorname{sinc}^2(0{,}1T_s f)\, e^{-j2\pi f\, 0{,}1T_s}\Bigr)\\
&= 0{,}01\,T_s^2\operatorname{sinc}^2(0{,}1T_s f)\, e^{-j2\pi f\, 0{,}1T_s}
\end{align*}
\begin{align*}
H_{eq}(f) &= \frac{T_s\, e^{j2\pi f\, 0{,}1T_s}}{0{,}01\,T_s^2\operatorname{sinc}^2(0{,}1T_s f)}\\
&= \begin{cases} \dfrac{2\cdot 10^{5}\, e^{j2\pi f\, 5\cdot 10^{-5}}}{\operatorname{sinc}^2(5\cdot 10^{-5} f)}, & |f| \le 1\ \mathrm{kHz}\ \text{ iken}\\[2ex] 0, & \text{else} \end{cases}
\end{align*}
\end{cozum}

\sayfa{4}
\sekilnotu[4]{$g(t)$ grafiği (dikey eksen $g(t)$, yatay eksen $t$): bir periyotluk sinüs benzeri eğri ($t=0$'dan başlayıp pozitif tepeye çıkar, sıfırdan geçer, negatif tepeye iner ve tekrar sıfıra döner). Eğri üzerinde turuncu, eşit aralıklı altı impuls oku: üçü pozitif (yukarı ok), üçü negatif (aşağı ok); okların boyu eğrinin değerini izliyor.}

\sekilnotu[4]{Aynı $g(t)$ eğrisi; bu kez aynı altı örnekleme anında turuncu dikdörtgen (düz tepeli) darbeler: üçü pozitif, üçü negatif.}

\sekilnotu[4]{$x_{PWM}$ grafiği: sabit yükseklikli, genişlikleri farklı dört dikdörtgen darbe (ilki $y$ ekseninde ve dar, sonra orta, en geniş ve tekrar daha dar); yatay eksen zaman.}

\sekilnotu[4]{$x_{PPM}(t)$ grafiği: ilk darbe $y$ ekseninde ve daha geniş; devamında sabit genişlikli dar darbeler, konumları kaydırılmış; referans konumlar kesikli dikey çizgilerle işaretli.}

$x_{PPM}(t)$: Pulse Width modulation\footnote{Claude notu: muhtemelen Pulse Position modulation olmalı.}

\sayfa{5}
\sekilnotu[5]{Bu sayfa basılı materyal içeriyor (ders kitabından şekil); dijitalleştirilmedi. Dört grafik üst üste: $g(t)$ (sinüs benzeri eğri, örnekleme okları), $x_{PAM}(t)$, $x_{PWM}(t)$ (basılı not: Pulse locations are the same but their widths change.) ve $x_{PPM}(t)$ (basılı not: Pulse widths are the same but their locations change.). El yazısı eklemeler: grafiklerin sağında PAM, PWM, PPM etiketleri ve altta aşağıdaki iki açıklama.}

PWM: Pulse Width Modulation (Darbe Genişlik Modülasyonu)

PPM: Pulse Position Modulation (Darbe Pozisyon Modülasyonu)

\sayfa{6}
\section{Time division multiplexing (TDM) \quad Zaman bölmeli çoklama (ZBÇ)}
\begin{itemize}
\item Multiple PAM signals are transmitted without overlapping.
\item At the receiver, PAM signals are separated by synchronization
\end{itemize}

\sekilnotu[6]{$x_a(t)$ grafiği: $\mathrm{BW}\{x_a(t)\}=B\ \mathrm{Hz}$; çan benzeri $x_a(t)$ eğrisi ve $0$, $T_s$, $2T_s$ anlarında turuncu dikdörtgen (PAM) darbeler.}

\sekilnotu[6]{$x_b(t)$ grafiği: $\mathrm{BW}\{x_b(t)\}=B\ \mathrm{Hz}$; azalan $x_b(t)$ eğrisi ve $T_s/2$ ile $3T_s/2$ anlarında kırmızı taralı dikdörtgen (PAM) darbeler.}

\sekilnotu[6]{Çoklayıcı: $x_1(t)$ girişi A noktasına ($[0,T_s/2]$ aralığı), $x_2(t)$ girişi B noktasına ($[T_s/2,T_s]$ aralığı) bağlı; dönen anahtar ($f_s$, not: $f_s\ge 2B$) çıkışı ``FTS ($\tau$)'' bloğuna verir; bloğa alttan ``TDM Clock, $f_p=2f_s$'' bloğu bağlı; çıkış $v(t)$. Altta $0$, $T_s/2$, $T_s$, $3T_s/2$, $2T_s$ anlarında impuls dizisi (``for N signals, $f_p=Nf_s$'').}

\sekilnotu[6]{$v(t)$ grafiği: PAM darbeleri sırayla $x_a(t)$ (turuncu) ve $x_b(t)$ (kırmızı) darbeleri olarak birbirini izliyor.}

\sekilnotu[6]{Ayırıcı: $v(t)$ girişi dönen anahtara ($f_s$) girer; anahtar A ve B kollarına bağlanır. A kolu: $x_a(t) \to [H_{eq}'(f)] \to x_1(t-T_0)$; B kolu: $x_b(t) \to [H_{eq}'(f)] \to x_2(t-T_0)$.}

FDM $\times$ TDM

Advantages: 1) Easier to implement \quad 2) Less crosstalk than FDM

Disadvantage: 1) Strict synchronization

Çok kaynaklı (Multi-source) $\to$ multiplexing (çoklama)

Çok kullanıcılı (Multi user) $\to$ Multiple Access (çoklu erişim)

\hspace*{1em}$\hookrightarrow$ TDMA (2.nesil, avrupada GSM), FDMA (1.nesil)

3.nesil: Code Division multiplexing

4.nesil: OFDM


\sayfa{7}
\section{Pulse code Modulation (Darbe kod modülasyonu - DKM)}

\elcizimi[0.9]{sekiller/Ders-5-s7-1.png}[7]

$x(t)$ analog signal \hfill $-x_p$ den $+x_p$ ye kadar

$x_p=\mathrm{Maks}\,|x(t)|$ \quad peak value (Tepe değer)

\underline{kuantalayıcı adımı} (Quantizer step) $=\Delta v=\dfrac{2x_p}{L}$ $\longrightarrow$ kuantalayıcıdaki seviye sayısı (\# of levels in quantizer)

\elcizimi[0.75]{sekiller/Ders-5-s7-2.png}[7]

Assume $\Delta v=1\ \mathrm{V}$

\begin{center}
\begin{tabular}{rcl}
$x(t)$ & & $x_q(t)$\\[3pt]
$1{,}07$ & $\longrightarrow$ & $1\ \mathrm{V}$\\
$1{,}40\ \mathrm{V}$ & $\longrightarrow$ & $1\ \mathrm{V}$\\
$0{,}55\ \mathrm{V}$ & $\longrightarrow$ & $1\ \mathrm{V}$\\
$2{,}25\ \mathrm{V}$ & $\longrightarrow$ & $2\ \mathrm{V}$\\
$-3{,}1\ \mathrm{V}$ & $\longrightarrow$ & $-3\ \mathrm{V}$
\end{tabular}
\end{center}

\sayfa{8}
\elcizimi[0.62]{sekiller/Ders-5-s8-1.png}[8]

8 seviyeyi kodluyoruz

$L=8$ \quad \textit{(Mid rise quantizer)}

Assume $\Delta v=1\ \mathrm{V}$

\begin{center}
\begin{tabular}{rcl}
\underline{$x(t)$} & & \underline{$x_q(t)$}\\[3pt]
$1{,}03\ \mathrm{V}$ & & $1{,}5\ \mathrm{V}$\\
$0{,}55\ \mathrm{V}$ & & $0{,}5\ \mathrm{V}$\\
$-3{,}1\ \mathrm{V}$ & & $-3{,}5\ \mathrm{V}$\\[6pt]
$1{,}01\ \mathrm{V}$ & $\longrightarrow$ & $1{,}5\ \mathrm{V}$
\end{tabular}
\qquad
\begin{tabular}{l}
$+\dfrac{\Delta v}{2}$ hata\\[4pt]
$-\dfrac{\Delta v}{2}$ hata
\end{tabular}
\end{center}

* Each sample is approximated with a quantized value, the approximation will result in an error less than $\Delta v/2$ in the positive direction or $-\Delta v/2$ in the negative \okunamadi{direction}

Quantization error : $(-\Delta v/2,\ \Delta v/2)$

\[
\underbrace{x_q(t)}_{\text{kuantalanmış işaret}} \quad = \quad \underbrace{x(t)}_{\text{Analog işaret}} + \overbrace{q(t)}^{\substack{\text{kuantalama gürültüsü}\\ \text{(Quantization noise)}}}
\]

$q$ is uniform in $[-\Delta v/2,\ \Delta v/2]$

Average quantization noise power $=N_q=E[q^2]$

(ort. kuantalama gürültüsü gücü)

\subsection{Reminder}
$x\sim U[a,b]$ düzgün dağ.

\elcizimi[0.3]{sekiller/Ders-5-s8-2.png}[8]

\[ \mu_x=\frac{a+b}{2} \qquad \sigma_x^2=\frac{(b-a)^2}{12} \]
\[ \sigma_q^2=E\!\left[(q-\mu_q)^2\right]=E[q^2]=N_q \]
\[ N_q=\frac{\left[\frac{\Delta v}{2}-\left(-\frac{\Delta v}{2}\right)\right]^2}{12}=\frac{\Delta v^2}{12} \]
\[ \Delta v=\frac{2x_p}{L}\ \Longrightarrow\ N_q=\frac{\Delta v^2}{12}=\frac{x_p^2}{3L^2} \]

İşaretin gücü $=S_0=\langle x^2(t)\rangle=P_x$

Gürültünün gücü $=N_0=N_q$

İşaretin gürültüye oranı (İGO)

Signal to noise Ratio (SNR)

\noindent\hfill signal to Quantization noise Ratio (SQNR) $\downarrow$

\formulkutu[önemli]{\mathrm{SNR}=\dfrac{S_0}{N_0}=\dfrac{P_x}{x_p^2}\cdot 3L^2}

\end{document}
